\section{Introducción}

Este trabajo\footnote[1]{Se utilizan los Lineamientos FACSO para la elaboración de este proyecto, donde los proyectos "(...) pueden seguir un formato de investigación empírica, discusión teórica, o intervención político social. El desarrollo del Seminario de Grado va a depender del tipo de enfoque que la o el estudiante decida adoptar. No obstante, hay elementos comunes que deberán ser abordados a lo largo de las secciones del Seminario de Grado. Estos incluyen antecedentes problematización, objetivos, perspectiva teórica y metodológica. El orden en que se presentan estos elementos puede variar (por ejemplo, la problematización puede situarse después de los antecedentes y perspectiva teórica)."} examina comparativamente la relación entre el ingreso familiar y el costo de vida en Chile, Brasil y Argentina durante el decenio 2015–2025. Este período se ha caracterizado por complejas transformaciones macroeconómicas, presiones inflacionarias disímiles o diferentes y reajustes estructurales en el mercado de trabajo regional \parencite{cepalPanoramaSocialAmerica2025}. Diversos estudios han evidenciado que la inflación no impacta de manera homogénea a la población, manifestando brechas diferenciadas según el nivel socioeconómico de los hogares \parencite{Adam2025, doValleGomesdaSilva2026} y forzando la adopción de estrategias de consumo de subsistencia ante el encarecimiento de bienes básicos \parencite{Nessieretal2026}. Asimismo, la dinámica entre la inflación de precios y la evolución salarial \parencite{Malheretal2026}, condicionada por la efectividad del salario mínimo \parencite{GindlingRonconi2025, MaurizioVazquez2016} y la heterogeneidad de los ingresos laborales \parencite{DuranKremerman2024, INE2024}, resulta determinante para evaluar las alteraciones en el poder adquisitivo y el grado de vulnerabilidad socioeconómica familiar.

A través de una perspectiva comparada y multidisciplinar, la investigación indaga en las mediaciones entre estas variaciones macroeconómicas y las condiciones materiales de vida de las familias en cada caso de estudio. 

Para ello, la propuesta se articula en un doble pilar analítico: por un lado, un cuerpo de antecedentes empíricos enfocado en la distribución del ingreso \parencite{ResumenResultadosCASEN2023}, la efectividad estabilizadora de las transferencias condicionadas y la protección social \parencite{Stampinietal2025, Carrere2026}, así como el rol de las instituciones laborales y el poder sindical \parencite{BarreraInsuaNoguera2026, DuranS2026}; por otro lado, un marco teórico que integra la economía política del trabajo y la acumulación de capital \parencite{DuranSilva2022, MarxElCapital2008} con los enfoques contemporáneos sobre estratificación social y reproducción del bienestar en América Latina \parencite{cepalPanoramaSocialAmerica2025, CostaRibeiroCarcalhaes2024}.